\subsection{条件溯源与可控匿名模型}

匿名机制能够保护用户隐私，但完全匿名也可能导致恶意行为难以追责。在外卖订单通信场景中，用户和骑手之间可能发生恶意骚扰、虚假投诉、配送纠纷或违规沟通等情况。如果系统只强调匿名而不保留管理能力，平台将难以处理异常事件；如果系统默认暴露真实身份，则又会削弱隐私保护效果。因此，本作品采用条件溯源与可控匿名相结合的设计思路。

\textbf{可控匿名}（Controllable Anonymity）是指系统在正常业务流程中隐藏用户真实身份，而在满足特定条件时，由具备权限的管理方按照审计流程恢复必要身份信息。该模型强调“日常匿名、异常可追责”。可控匿名区别于完全匿名（如 Tor）和完全实名两种极端模式，在隐私保护与责任追踪之间寻求平衡。

在 FoodShield 系统中，普通通信状态下，骑手端无法看到用户真实身份，也无法看到用户 PID；管理员端默认只查看订单摘要、消息哈希、Merkle Root 和完整性验证结果。当发生投诉、纠纷或平台审计需求时，用户或骑手可以基于订单号提交溯源申请，管理员经过身份认证和权限校验后，才能触发条件溯源流程。

条件溯源流程可以抽象为：

\begin{equation}
\mathit{orderID} \xrightarrow{\operatorname{verify\text{-}integrity}} \mathit{PID} \xrightarrow{\operatorname{lookup}} \mathit{userID}
\end{equation}

其中，$\mathit{orderID}$是业务入口，$\mathit{PID}$是系统内部匿名标识，$\mathit{userID}$是真实用户身份标识。系统不鼓励管理员直接通过 PID 发起溯源，而是以订单号作为溯源入口。这种方式更符合外卖平台的真实业务逻辑（投诉和纠纷通常以订单为单位发起），也可以降低 PID 被滥用或过度暴露的风险。

条件溯源流程包含两个关键安全约束：

\begin{itemize}[leftmargin=2em]
\item \textbf{完整性前置校验}：在执行溯源之前，系统必须先对目标订单的通信日志进行 Merkle Tree 完整性验证。只有当日志完整性验证通过时，才允许继续溯源流程。这一约束确保了溯源所依据的通信记录未被篡改，防止攻击者通过篡改日志来误导溯源结果。
\item \textbf{关键词条件触发}：溯源过程通常由投诉或举报触发，系统根据投诉内容中的关键词进行命中匹配，确认是否存在异常通信行为。只有当关键词命中且完整性校验通过时，才执行从$\mathit{orderID}$到$\mathit{userID}$的身份映射。这一双重门槛机制防止管理员滥用溯源权限进行无依据的身份查询。
\end{itemize}

在条件溯源过程中，系统应记录管理员的关键操作，包括登录行为、查看订单摘要、执行完整性验证、确认溯源申请和查询身份映射等。每一次溯源操作都应写入审计日志，记录操作者、操作时间、目标订单、溯源理由和操作结果。通过这种方式，系统不仅可以追踪用户或骑手的异常行为，也可以约束管理员权限，避免审计能力本身成为新的隐私风险。

% !TEX root = ../../main.tex
% 条件溯源概念流程图（第二章，简化版主干流程）
\begin{figure}[H]
\centering
\begin{tikzpicture}[
  node distance = 0.55cm and 1.6cm,
  proc/.style   = {rectangle, draw, rounded corners=3pt, minimum width=3.0cm, minimum height=0.65cm,
                   font=\small, fill=blue!8, draw=blue!50!black},
  decision/.style = {diamond, draw, aspect=2.2, minimum width=3.2cm, minimum height=0.8cm,
                     font=\small, align=center, fill=orange!15, draw=orange!60!black, inner sep=2pt},
  term/.style   = {rectangle, draw, rounded corners=8pt, minimum width=3.0cm, minimum height=0.65cm,
                   font=\small, fill=green!12, draw=green!50!black},
  reject/.style = {rectangle, draw, rounded corners=8pt, minimum width=2.2cm, minimum height=0.62cm,
                   font=\small, fill=red!12, draw=red!50!black},
  arr/.style    = {-Stealth, thick},
  lbl/.style    = {font=\small},
]

\node[proc]   (start)  {投诉申请（orderID + 原因）};
\node[decision, below=0.8cm of start] (d1) {完整性\\校验通过？};
\node[reject,  right=2.2cm of d1]  (r1) {拒绝溯源（日志可疑）};
\node[decision, below=0.9cm of d1] (d2) {关键词\\命中？};
\node[reject,  right=2.2cm of d2]  (r2) {拒绝溯源（无据可查）};
\node[proc,    below=0.9cm of d2]  (pid) {PID 溯源};
\node[term,    below=0.7cm of pid] (uid) {获得 userID（真实身份）};
\node[proc,    below=0.6cm of uid] (log) {写入审计日志};

\draw[arr] (start) -- (d1);
\draw[arr] (d1) -- node[lbl, left]{是} (d2);
\draw[arr] (d1) -- node[lbl, above]{否} (r1);
\draw[arr] (d2) -- node[lbl, left]{是} (pid);
\draw[arr] (d2) -- node[lbl, above]{否} (r2);
\draw[arr] (pid) -- (uid);
\draw[arr] (uid) -- (log);

\end{tikzpicture}
\caption{条件溯源概念流程（预备知识层）}
\label{fig:trace_concept}
\end{figure}


综上，条件溯源与可控匿名模型是 FoodShield 系统区别于普通匿名通信系统的重要设计。它既避免了普通外卖聊天系统中过度暴露身份的问题，也避免了完全匿名系统中责任主体无法追踪的问题，从而在隐私保护、订单认证和平台管理之间形成平衡。
